%% CV for Continia Software - Commercial Enablement Specialist (Dansk)
%% Compiled with: lualatex

\documentclass[11pt,a4paper,sans]{moderncv}
\moderncvstyle{banking}
\moderncvcolor{blue}

\renewcommand*{\firstnamestyle}[1]{{\fontsize{34}{36}\bfseries\upshape\color{color1}#1}}
\renewcommand*{\lastnamestyle}[1]{{\fontsize{34}{36}\bfseries\upshape\color{color1}#1}}
\renewcommand*{\sectionstyle}[1]{{\sectionfont\color{color1}#1}}

\usepackage[utf8]{inputenc}
\usepackage{hyperref}
\hypersetup{
    colorlinks=true,
    linkcolor=blue,
    filecolor=magenta,
    urlcolor=blue,
    pdftitle={Jacob El-Omar - CV},
    pdfpagemode=FullScreen,
}
\usepackage[scale=0.80]{geometry}
\usepackage{import}

\name{Jacob}{El-Omar}
\address{Aalborg, Danmark}{}{}
\phone[mobile]{+45 53 35 27 82}
\email{elomarjc@gmail.com}
\extrainfo{\href{https://linkedin.com/in/jacob-el-omar}{LinkedIn} \textbullet{} \href{https://github.com/elomarjc}{GitHub}}

\begin{document}

\makecvtitle

% ============================================================
%     PROFIL
% ============================================================

\vspace{6pt}
\small{Nyuddannet civilingeniør i elektroniske systemer fra Aalborg Universitet med stærke tekniske og strukturelle kompetencer samt erfaring med IT-platforme, datamodellering og procesoptimering. Erfaren i at omsætte komplekse tekniske koncepter til strukturerede workflows og præsentationsmaterialer. Anvender aktivt AI-kodningsagenter som Claude Code til automatisering af opgaver og dokumentation. Motiveret for at støtte kommercielle strategier og salgsprocesser hos Continia Software i Nørresundby.}

% ============================================================
%     FAGLIGE KOMPETENCER
% ============================================================

\section{Faglige Kompetencer}
\vspace{1pt}
\begin{itemize}

\item \textbf{Struktur \& Procesoptimering}: Omsætning af komplekse tekniske og organisatoriske strategier til overskuelige strukturer og workflows.

\item \textbf{IT- og Platformudvikling}: Kendskab til softwarearkitektur, databaser (PostgreSQL, MySQL), C\#, Python og REST-API'er.

\item \textbf{AI-Understøttet Effektivisering}: Aktiv brug af agentiske AI-værktøjer (Claude Code) til hurtig opgavestrukturering, dokumentation og indholdsproduktion.

\item \textbf{Formidling \& Brugerdialog}: Trænet i problembaseret læring (PBL) med fokus på tæt dialog, vidensdeling og præsentation af tekniske løsninger.

\end{itemize}

% ============================================================
%     ERHVERVSERFARING
% ============================================================

\section{Erhvervserfaring}
\vspace{3pt}
\begin{itemize}

\item{\cventry{Jan 2026--Nu}{IT- \& Webudvikler (Frivillig)}{Dawah Danmark}{Aalborg, Danmark}{}{\vspace{1pt}
\begin{itemize}
    \item Moderniserede og strukturerede organisatorisk IT-platform og databaser på WordPress.
    \item Optimerede brugerflade og administrative processer for bedre intern brugbarhed.
\end{itemize}}}

\vspace{3pt}

\item{\cventry{Sep 2024--Jun 2025}{AI/ML Engineer (Praktik)}{Ai Health Highway}{Remote}{}{\vspace{1pt}
\begin{itemize}
    \item Udviklede og validerede støjfiltreringsalgoritmer i Python til medicinske enheder.
    \item Strukturerede teknisk dokumentation og afrapporterede resultater til internationale interessenter.
\end{itemize}}}

\vspace{3pt}

\item{\cventry{Jun 2023--Aug 2024}{Lager- og Logistikmedarbejder}{Lyn Express}{Aalborg, Danmark}{}{\vspace{1pt}
\begin{itemize}
    \item Styrede sortering og logistik af forsendelser i et tempofyldt miljø under stramme deadlines.
    \item Automatiserede administrative rapporteringsrutiner i Excel for øget effektivitet.
\end{itemize}}}

\end{itemize}

% ============================================================
%     UDDANNELSE
% ============================================================

\section{Uddannelse}
\vspace{1pt}
\begin{itemize}

\item{\cventry{Sep 2023--Jun 2025}{Civilingeniør, cand.polyt. (Elektroniske Systemer)}{Aalborg Universitet}{Aalborg, Danmark}{}{\vspace{1pt}
Speciale: ``Adaptive Noise Cancellation For Electronic Stethoscopes.'' Integrerede algoritmisk datamodellering og støjfiltrering.
}}

\vspace{3pt}

\item{\cventry{Sep 2020--Aug 2023}{Bachelor i teknisk videnskab (Elektronik og IT)}{Aalborg Universitet}{Aalborg, Danmark}{}{\vspace{1pt}
Specialisering i reguleringsteknik og indlejrede systemer.
}}

\end{itemize}

% ============================================================
%     UDVALGTE PROJEKTER
% ============================================================

\section{Udvalgte Projekter}
\vspace{1pt}
\begin{itemize}

\item{\textbf{Grow a Fish (Fritidsprojekt)}: Egenudviklet spil i Flutter/Dart med en Supabase PostgreSQL-backend og REST-API'er.}

\item{\textbf{Adaptiv Signalfiltrering (Speciale, Universitetsprojekt)}: Udviklede en databehandlingspipeline i Python til statistisk evalueringsmodel.}

\item{\textbf{Embedded Faldsikringssystem (Universitetsprojekt)}: Programmerede et realtids fall-detection system på Cypress PSoC i C.}

\end{itemize}

% ============================================================
%     SPROG
% ============================================================

\section{Sprog}
\vspace{1pt}
\begin{itemize}
\item Dansk (modersmål), Engelsk (flydende).
\end{itemize}

% ============================================================
%     REFERENCER
% ============================================================

\section{Referencer}
\vspace{1pt}
\begin{itemize}
\item Afleveres efter aftale.
\end{itemize}

\end{document}
